% methods.tex (국문 v13): 자료와 방법(결과 앞). 번호 없는 소절, 수식 (1)-(9).
\section{자료와 방법}\label{sec:methods}

\subsection{연구 지역과 활동층 관측}
여러 해에 걸친 융해 깊이 관측이 조밀한 세 지역, 알래스카 · 레나 델타 · 캐나다 북서부(그림과 표에서는 Canada (NW))를 대상으로 하였다(그림~\ref{fig:data}, 표~\ref{tab:data}). 알래스카와 캐나다는 ABoVE 토양 수분 · 활동층 두께 자료 v2\cite{Moore2025}와 CALM 계절 말 격자 값\cite{Streletskiy2025}을, 레나 델타는 ALLena 융해 깊이 자료\cite{Veremeeva2025}를 사용하였다. ABoVE 탐침 측정은 8--10월, $0<\mathrm{ALT}<300$ cm인 값으로 제한하였다. 그보다 이른 측정은 계절 최대값이 아니라 그 시점의 융해 깊이를 나타내기 때문이다. ALLena는 탐침 한계 표지가 없는 8--9월 값으로 제한하고, 위치 · 연도별 계절 최댓값을 취하였다. 관측은 좌표를 0.0001$^\circ$(약 10 m)로 반올림한 위치인 지점별 · 연도별로 묶어 평균하고 평균 관측일을 보존하였다. 따라서 ABoVE 횡단선의 탐침 지점은 각각 별개의 지점이며, 한 1 km 셀 안에 여러 지점이 있을 수 있다. 날짜가 없는 CALM 값은 244번째 날(9월 1일)로 두었다. 연도별 기후 입력이 있는 2010--2024년만 사용하였다. 라벨은 알래스카 \nncyAk{}개 지점-연도(\nyrsAkko), 레나 델타 \nncyLe{}개(\nyrsLeko), 캐나다 \nncyCa{}개(\nyrsCako)이며, 0.5$^\circ$ 블록 \nnblkAk, \nnblkLe, \nnblkCa{}개에 분포한다(표~\ref{tab:data}). 지점은 1 km 셀 \nnkmcellAk{}개, \nnkmcellLe{}개, \nnkmcellCa{}개(셀-연도 \nnkmcyAk{}개, \nnkmcyLe{}개, \nnkmcyCa{}개)에 들어 있고, 한 횡단선의 지점이 한 셀 · 한 해에 최대 \nnkmmaxAk{}개까지 몰리므로 지점-연도는 1 km 규모에서 독립 표본이 아니며 공간 블록 설계가 이를 다룬다. 자료 출처 · 판 · 약관은 보충 표 S1, 분석별 사용 연도는 보충 표 S2, 처리 단계는 보충 방법 S1에 정리하였다.

% table1.tex (국문 v13): 표 1. 지역별 라벨과 입력 · 제품의 기간 · 사용 연도. 수치는 매크로.
\begin{table}[!htbp]
\caption{자료와 사용 연도. 위: 지역별 ALT 라벨(관측 연도 기후 입력이 있는 2010--2024년 지점-연도; 지점은 좌표를 0.0001$^\circ$ 로 반올림한 관측 위치이고 1 km 셀은 지점을 담는 지도 셀이다). 아래: 입력 자료와 제품의 제공 기간과 이 연구에서 쓴 연도. 블록은 공간 교차검증과 부트스트랩에 쓰는 0.5$^\circ$ 칸이며, 캐나다는 2020년 라벨이 없다.}\label{tab:data}
\scriptsize
\begin{tabular}{@{}lrrrrlrl@{}}
\toprule
지역 & 지점-연도 & 지점 & 1 km 셀 & 블록 & 연도 & 4개 해 이상 지점 & 최다 연도 \\
\midrule
알래스카 & \nncyAk{} & \nncellAk{} & \nnkmcellAk{} & \nnblkAk{} & 2010--2024 & \nnmultiAk{} & \ntopyrAk{} (\ntopyrshareAk\%) \\
레나 델타 & \nncyLe{} & \nncellLe{} & \nnkmcellLe{} & \nnblkLe{} & 2010--2022 & \nnmultiLe{} & \ntopyrLe{} (\ntopyrshareLe\%) \\
캐나다 북서부 & \nncyCa{} & \nncellCa{} & \nnkmcellCa{} & \nnblkCa{} & 2010--2024 & \nnmultiCa{} & \ntopyrCa{} (\ntopyrshareCa\%) \\
외부 CALM 지역 & \nftScoredRows{} & \nftScoredCells{} & -- & \nftScoredBlocks{} & \nyrsext{} & -- & -- \\
\bottomrule
\end{tabular}

\vspace{6pt}
\begin{tabular}{@{}>{\raggedright\arraybackslash}p{0.26\textwidth}>{\raggedright\arraybackslash}p{0.11\textwidth}>{\raggedright\arraybackslash}p{0.13\textwidth}>{\raggedright\arraybackslash}p{0.16\textwidth}>{\raggedright\arraybackslash}p{0.16\textwidth}@{}}
\toprule
자료 & 해상도 & 제공 기간 & 사용 연도 & 역할 \\
\midrule
ERA5-Land 월평균\cite{munozsabater2021_era5_land} & 0.1$^\circ$ & 1950년부터 & 2010--2024년, 평균 2011--2024년 & 기후 입력 \\
Copernicus DEM\cite{CopDEM2021} & 30 m & 정적 & -- & 지형 \\
SoilGrids 2.0\cite{Poggio2021} & 250 m & 정적 & -- & 토양 \\
MODIS MOD13Q1, MOD10A2\cite{Didan2021,HallRiggs2021} & 250--500 m & 2000년부터 & 2010--2024년 & 민감도 시험 \\
ESA CCI 영구동토 확률\cite{Westermann2024} & 1 km & 1997--2021년 & 평균 & 지도 범위 \\
Wei 등\cite{Wei2026} & 1 km & 2000--2024년 & \nyrsprodWei{}년 & 비교 \\
Liu 등\cite{Liu2024} & 1 km & 2003--2020년 & \nyrsprodLiu{}년 & 비교 \\
ESA CCI ALT v4\cite{Westermann2024} & 1 km & 1997--2021년 & \nyrsprodCcifour{}년 & 비교 \\
ESA CCI ALT v5\cite{Westermann2025} & 1 km & 1997--2023년 & \nyrsprodCcifive{}년 & 비교 \\
Yi와 Kimball\cite{YiKimball2019} & 1 km & 2001--2015년 & \nyrsprodYk{}년(알래스카) & 비교 \\
Aalto 등\cite{Aalto2018} & 30$''$ & 2000--2014년 & 시간 불변 & 비교 \\
\bottomrule
\end{tabular}
\end{table}


\paperfig{Fig1_data}{1.0}{\textbf{연구 지역, 관측, 자료 기간.} \textbf{a}, ESA CCI 영구동토 확률(1997--2021년 평균)과 세 연구 지역, 2010--2024년 ALT 라벨이 있는 관측 지점 수. \textbf{b}, 관측 연도별 지점-연도 라벨 수(2010--2024년). \textbf{c}, 라벨, 기후 입력에 쓴 ERA5-Land 월자료, 정적 공변량, 공개 제품 6종의 기간. 제품마다 제공 기간과 라벨 셀에서 비교에 쓴 연도를 구분하였고, 음영은 분석 기간 2010--2024년이다.}{fig:data}

\subsection{관측 연도 기후 입력}
2010년 1월--2024년 12월 ERA5-Land 월자료\cite{munozsabater2021_era5_land}(2 m 기온, 표층 토양 온도, 적설 수당량)를 각 지점의 가장 가까운 육지 격자에서 추출하였다. 셀 $c$, 연도 $y$의 융해 도일은 다음과 같다.
\begin{equation}
\mathrm{TDD}_{c,y} = \sum_{m=1}^{12} \max\!\left(T_{c,y,m},\,0\right) d_{y,m}
\label{eq:tdd}
\end{equation}
여기서 $T_{c,y,m}$은 월평균 기온, $d_{y,m}$은 $m$월의 일수이다. 관측은 대부분 융해 계절이 끝나기 전에 이루어졌으므로, 모델 입력은 관측일 $t$까지 누적한 TDD이다.
\begin{equation}
\mathrm{TDD}^{\mathrm{obs}}_{c,y} = C_{c,y,k-1} + \frac{t - t_{k-1}}{d_{y,k}} \left( C_{c,y,k} - C_{c,y,k-1} \right)
\label{eq:tddobs}
\end{equation}
$C_{c,y,k}$는 $k$월 말까지의 누적 TDD, $k$는 $t$가 속한 달, $t_{k-1}$은 전달의 마지막 날이다. 이와 함께 기후 변수 일곱 개, 즉 연평균 기온, $\sqrt{\mathrm{TDD}}$, 가장 따뜻한 달과 가장 추운 달의 평균 기온, 직전 겨울(10월 $y-1$--4월 $y$)의 동결 도일, 표층 토양 온도 평균, 수문년(10월 $y-1$--9월 $y$) 적설 수당량 평균을 셀 · 연도별로 계산하였다(보충 표 S3, 보충 방법 S2). 같은 지점-연도에서 두 가지 대안 입력을 함께 평가하였다. 관측 연도의 계절 전체 TDD와, 연도별 값의 2011--2024년 평균(평년 기후 입력)이다. 후자는 기존 제품의 다년 입력에 해당한다. 연도별 지도는 CALM 프로토콜의 계절 말 기준일인 9월 1일(244일)까지 누적한 TDD를 쓴다. 이 입력은 학습 라벨의 범위 안(관측일 사분위 범위 224--243일)에 있으며, 계절 전체 합은 이를 20--30\% 넘는다. 따라서 연도별 지도 값은 9월 1일의 ALT를 나타낸다.

\subsection{정적 공변량}
지형 변수 여섯 개(고도, 경사, 사면 방향의 사인 · 코사인, 지형 위치 지수, 거칠기)는 30 m Copernicus DEM\cite{CopDEM2021}에서 각 셀 둘레 $33\times33$ 화소 창(약 1 km)으로 계산하였다. 토양 변수 아홉 개(5--15 cm의 점토, 모래, 실트, 용적 밀도, 조립질, pH와 0--5, 5--15, 15--30 cm의 토양 유기탄소)는 250 m SoilGrids 2.0\cite{Poggio2021}에서 추출하였다. 따라서 주 특징 벡터 $\mathbf{x}$는 기후 8개, 지형 6개, 토양 9개의 23개 변수로 이루어진다. 민감도 시험에는 기후만, 기후와 지형, 주 묶음에 2010--2024년 MODIS 여름 NDVI · EVI와 눈 소실 시기(MOD13Q1 · MOD10A2 v6.1)\cite{Didan2021,HallRiggs2021}를 더한 묶음, 주 묶음에 고해상 토양 탄소 · 피복 분류\cite{Palmtag2022}를 더한 묶음, 전년 기후를 포함한 전체 변수를 사용하였다. ESA CCI 영구동토 확률(1997--2021년 평균)\cite{Westermann2024}은 예측 변수가 아니라 지도 범위(확률 50\% 이상)를 정하는 데에만 사용하였다.

\subsection{모델 구성}
Stefan 모델은 융해 깊이를 다음과 같이 나타낸다.
\begin{equation}
\widehat{\mathrm{ALT}}_{\mathrm{S}} = E\,\sqrt{\mathrm{TDD}^{\mathrm{obs}}}
\label{eq:stefan}
\end{equation}
$E$(cm\,($^\circ$C\,d)$^{-1/2}$)는 토양의 열적 성질과 수분을 요약한다\cite{Stefan1891,Riseborough2008}. 이 관계는 레나 델타 관측 지점에서 계절 안에 잘 성립하며, 반복 관측된 셀에서는 연도 간에도 더 큰 산포와 함께 성립한다(보충 그림 S1). $n$개 학습 라벨 $y_i$와 $s_i = \sqrt{\mathrm{TDD}^{\mathrm{obs}}_i}$에 대해 계수는 원점을 지나는 최소제곱으로 재보정한다.
\begin{equation}
\hat{E} = \frac{\sum_{i=1}^{n} y_i\, s_i}{\sum_{i=1}^{n} s_i^{2}}
\label{eq:E}
\end{equation}
원천 계수 Stefan 모델은 다른 두 지역에서 맞춘 $\hat{E}$를, 재보정 Stefan 모델은 대상 지역 학습 라벨로 맞춘 $\hat{E}$를 쓴다. 물리 잔차 모델은 재보정 Stefan 예측을 기준값으로 두고 학습한 보정량을 더한다(그림~\ref{fig:workflow}).
\begin{equation}
\widehat{\mathrm{ALT}} = \hat{E}\, s + \lambda\, g(\mathbf{x}), \qquad g = \arg\min_{f} \sum_{i=1}^{n} \left[ \left( y_i - \hat{E}\, s_i \right) - f(\mathbf{x}_i) \right]^{2}
\label{eq:residual}
\end{equation}
$g$는 CatBoost gradient boosting 회귀기\cite{Prokhorenkova2018}(반복 600, 학습률 0.03, 깊이 6, L2 정칙화 3)이며, $\lambda=0.5$는 보정량을 물리 기준값 쪽으로 수축한다($\lambda$ = 0.25, 1과 내부 블록 교차검증으로 고른 값을 민감도로 시험하였다). 비교한 세 구조는 물리 정보를 다른 위치에 넣는다. 직접 ML은 $\mathbf{x}$만으로 ALT를 학습하고, 물리 입력 ML은 재보정 Stefan과 Kudryavtsev\cite{Kudryavtsev1974} 예측을 특징으로 더하며, 물리 유사라벨은 채점 겹 셀의 공변량과 학습 겹 Stefan 예측으로 만든 같은 수의 행을 학습 자료에 더한다. 모든 학습 모델은 무작위 씨앗 3개로 학습해 예측을 평균하였다. 학습기 민감도에서는 CatBoost 대신 랜덤 포레스트, 다층 퍼셉트론, TabM, FT-Transformer\cite{Gorishniy2021,Gorishniy2025,Holzmuller2024}를 사용하였고, 지역별 전체 라벨로 학습한 잔차 모델과 직접 모델의 SHAP 값\cite{Lundberg2017}을 계산하였다.

\paperfig{Fig2_workflow}{1.0}{\textbf{지점-연도 설계, Stefan 기준값, 물리 잔차 모델, 검증 설계.} \textbf{a}, 각 라벨은 한 해의 관측 지점 하나이다. 왼쪽은 관측일의 융해 깊이와 계절 말의 ALT, 가운데는 레나 델타 셀의 2013년 누적 융해 도일과 그 셀의 2011--2024년 평균, 다른 해이다. 모델 입력은 관측일까지 누적한 값(TDD$_{\mathrm{obs}}$)이고 연도별 지도는 9월 1일까지 누적한 값을 쓴다. 오른쪽은 알래스카 지점-연도 \nncyAk{}개의 관측 ALT 와 $\sqrt{\mathrm{TDD_{obs}}}$ 의 관계이며, 원점을 지나는 최소제곱으로 적합한 Stefan 기준값(식 (\ref{eq:E}))과 잔차 5개를 표시하였다. \textbf{b}, 물리 잔차 모델: Stefan 기준값과, 기후 · 지형 · 토양 입력 23개로 그 잔차를 학습하는 gradient boosting 을 가중치 $\lambda = 0.5$ 로 더한다(식 (\ref{eq:residual})). 지도는 2014년 9월 1일까지 누적한 TDD 를 쓰며(알래스카), 입력 썸네일은 변수마다 2--98 백분위로 늘린 고유 색 눈금을 쓴다. \textbf{c}, 레나 델타 라벨 셀의 검증 설계(모든 겹 표시): 점 색은 겹(지점-연도 단위 무작위 분할, 0.5$^\circ$ 공간 블록은 겹 색 테두리), 관측 연도(연도 홀드아웃 · 블록 × 연도, 회색은 지점-연도 20개 미만으로 채점하지 않은 해), 시험 지역(지역 홀드아웃, 모든 셀)을 뜻한다. 오른쪽은 2011--2024년 평균 대비 알래스카 지도의 2012년 · 2019년 ALT 편차.}{fig:workflow}

\subsection{검증 설계와 불확실성}
모든 모델은 다섯 가지 설계의 겹 밖 예측으로 평가하였다\cite{Roberts2017,Ploton2020}. 지점-연도 단위 무작위 5겹 교차검증, 0.5$^\circ$ 블록을 지점-연도 수가 고르게 겹에 배정한 공간 블록 5겹 교차검증(주 설계), 지점-연도 20개 이상인 각 관측 연도를 나머지 해로 학습한 모델로 예측하는 연도 홀드아웃, 채점 지점-연도의 블록 겹과 연도를 모두 학습에서 빼는 블록 × 연도 홀드아웃, 시험 지역에서 100 km 안의 학습 셀을 뺀 뒤 다른 두 지역으로 학습한 모델로 예측하는 지역 홀드아웃이다. 정확도는 다음과 같이 나타낸다.
\begin{equation}
\mathrm{RMSE} = \left[ \frac{1}{n} \sum_{i=1}^{n} \left( \hat{y}_i - y_i \right)^{2} \right]^{1/2}, \qquad \Delta\mathrm{RMSE}_{A,B} = \mathrm{RMSE}_{A} - \mathrm{RMSE}_{B}
\label{eq:rmse}
\end{equation}
$\Delta\mathrm{RMSE}$의 두 항은 같은 지점-연도에서 계산하며, 음수이면 방법 $A$가 낫다. 불확실성은 짝지은 블록 부트스트랩으로 정량화하였다. 0.5$^\circ$ 블록을 복원 추출로 $B=2000$회(라벨 수 실험은 1000회) 재추출하고, 재추출된 블록의 모든 지점-연도로 $\Delta\mathrm{RMSE}$를 다시 계산하여, 반복값의 2.5와 97.5 백분위수를 95\% 구간으로 정하였다.
\begin{equation}
\mathrm{CI}_{95} = \left[ Q_{0.025}\!\left( \Delta\mathrm{RMSE}^{(b)} \right),\; Q_{0.975}\!\left( \Delta\mathrm{RMSE}^{(b)} \right) \right], \quad b = 1,\ldots,B
\label{eq:ci}
\end{equation}
구간이 0을 제외하면 차이가 판정된 것으로 기술한다. 이득의 조건은 공변량 여덟 개(관측 연도의 계절 전체 TDD, 동결 도일, 적설 수당량, 토양 유기탄소, 점토, 경사, 고도, 여름 NDVI)의 지역 안 3분위에서 분석하였고, 지점-연도 30개 또는 블록 5개 미만인 3분위는 평가하지 않았다. 연도별 RMSE는 지점-연도 30개 이상인 해에서, 관측과 예측의 연도 간 변화 상관은 여러 해 관측된 셀에서 계산하였다.

\subsection{라벨 수 · 전이 실험}
지역 안 실험에서는 공간 블록 겹마다 학습 라벨을 $n$ = 25, 50, 100, 250, 500, 1000, 2500개 또는 전량으로 줄이고(무작위 3회), 같은 $n$개 라벨로 Stefan 계수를 재보정한 뒤 시험 겹 전체를 채점하였다. 전이 실험에서는 다른 두 지역의 전체 라벨과 대상 지역 학습 겹에서 뽑은 $n$ = 0, 25, 50, 100, 250, 500개 라벨로 학습한 모델로 대상 지역의 공간 블록 겹을 채점하였다. $\hat{E}$는 대상 라벨 $n$개로 다시 맞추었고($n\geq3$), 잔차 학습기는 원천 계수 기준값을 쓰는 원천 행과 다시 맞춘 기준값을 쓰는 대상 행으로 학습하였다. 외부 시험에는 세 연구 지역과 산악 영구동토(스위스, 카자흐스탄) 밖의 2010--2024년 CALM 관측\cite{Streletskiy2025}을 사용하였다. 검열값, 이른 측정 · 화재 · 침수 · 영구동토면 미도달 표지가 있는 값, 연구 라벨에서 100 km 안의 지점을 제외하여 8개 지역의 \nextcells{}개 셀, \nextcy{}개 지점-연도를 남겼다(보충 방법 S3). 각 외부 지역은 블록 단위로(블록이 하나인 그린란드는 셀 단위로) 채점하였으며, 세 연구 지역과 같은 외부 지역의 다른 블록에서 뽑은 라벨 $n$ = 0, 5, 10, 25, 50개 또는 전량으로 학습하였다. 학습 범위 밖인 평균 ALT 약 6 m의 스발바르 시추공 한 곳은 채점에서 제외하여, 지점-연도 \nftScoredRows{}개, 셀 \nftScoredCells{}개, 블록 \nftScoredBlocks{}개가 남았다. 실행 목록은 보충 표 S4에 있다.

\subsection{공개 제품 비교}
라벨 셀에서 제품 6종을 추출하였다. Wei 등\cite{Wei2026}(1 km, 2000--2024년), Liu 등\cite{Liu2024}(1 km, 2003--2020년), ESA CCI Permafrost ALT v4(1997--2021년)와 v5(1997--2023년)\cite{Westermann2024,Westermann2025}, Yi와 Kimball\cite{YiKimball2019}(알래스카, 2001--2015년), Aalto 등\cite{Aalto2018}(2000--2014년 기준기, 셀당 값 하나)이다. 연도별 제품은 관측 연도의 값을 사용하였다(보충 방법 S4). 제품과 모델은 Wei 등, Liu 등, ESA CCI v4가 모두 관측 연도 값을 가진 공통 지점-연도(\nyrscc{}년)에서 비교하였다. Wei 등은 학습 지점 좌표를 공개하였으며, 누설 마스크는 이 지점에서 25 km 안의 라벨 셀을 제외한다. Liu 등은 이 연구에서 쓴 ABoVE 관측의 이전 판으로 학습하였으므로\cite{Liu2024}, 알래스카 · 캐나다 값은 우리 라벨과 독립이 아니다. 잔차 학습기가 제품도 보정할 수 있는지 시험하기 위해, 관측 연도의 제품 값 $p$를 학습 겹에서 다음과 같이 재보정하고
\begin{equation}
\hat{p} = a + b\, p
\label{eq:prod}
\end{equation}
잔차 $y-\hat{p}$를 식~(\ref{eq:residual})과 같은 공변량, 학습기, $\lambda$, 겹으로 학습하였다(보충 방법 S5).

\subsection{연도별 지도, 적용 가능 영역, 예측 구간}
세 지역 각각에서 재보정 Stefan 모델과 물리 잔차 모델을 지역의 모든 지점-연도로 학습하고, ESA CCI 영구동토 확률 50\% 이상인 1 km 셀(\nncellmapAk{}개, \nncellmapLe{}개, \nncellmapCa{}개)에 2011--2024년 각 해의 기후(9월 1일까지 누적한 TDD)로 적용하였다. 적용 가능 영역은 Meyer와 Pebesma\cite{MeyerPebesma2021}를 따랐다. $\mathbf{z}$를 표준화한 23차원 특징 벡터(2011--2024년 평균 기후)라 할 때, 지도 셀 $k$의 비유사도는 모든 학습 행 $j$에 대한 $d_k = \min_{j} \lVert \mathbf{z}_k - \mathbf{z}_j \rVert_2$이며, 다음을 만족하면 적용 가능 영역 안으로 본다.
\begin{equation}
d_k \leq \tau, \qquad \tau = Q_{0.95}\!\left\{ \min_{j:\, b(j) \neq b(i)} \lVert \mathbf{z}_i - \mathbf{z}_j \rVert_2 \;:\; i \in \text{학습 행} \right\}
\label{eq:aoa}
\end{equation}
$b(\cdot)$는 공간 블록이다. 지도 값의 90\% 예측 구간은 공간 블록 겹 밖 잔차의 5와 95 백분위수로 정하였다. 표시용으로 1 km 값을 250 m 음영 격자에 쌍선형 보간하였다.

\subsection{계산}
모델은 씨앗을 고정해 클라우드 노드(Rescale; 96코어 CPU 노드, 신경망 학습기는 NVIDIA A100 GPU 노드 1대)에서 학습하였고, 자료 준비 · 집계 · 그림은 로컬 서버에서 수행하였다(보충 표 S4).

\subsection{코드 가용성}
자료 준비, 모델링, 평가, 그림 코드는 [저장소 · DOI]에서 받을 수 있다. 본문의 모든 수치는 원고 매크로를 만드는 스크립트가 저장된 결과 표에서 생성한다.
